\documentclass[11pt]{article}

\usepackage[margin=1in]{geometry}
\usepackage{amsmath,amssymb,amsthm}
\usepackage{enumitem}
\usepackage[hidelinks]{hyperref}
\urlstyle{same}

\newtheorem{theorem}{Theorem}
\newtheorem{proposition}[theorem]{Proposition}
\newtheorem{corollary}[theorem]{Corollary}
\newtheorem{lemma}[theorem]{Lemma}
\theoremstyle{definition}
\newtheorem{definition}[theorem]{Definition}
\theoremstyle{remark}
\newtheorem{remark}[theorem]{Remark}
\newtheorem{question}[theorem]{Question}

\newcommand{\F}{\mathbb{F}}
\newcommand{\wt}{\operatorname{wt}}
\newcommand{\supp}{\operatorname{supp}}
\newcommand{\Span}{\operatorname{span}}
\newcommand{\taglabel}[1]{\textnormal{\textbf{[#1]}}}
\newcommand{\PROVED}{\taglabel{PROVED}}
\newcommand{\STANDARD}{\taglabel{STANDARD MATH}}
\newcommand{\UNREVIEWED}{\taglabel{AWAITING ADVERSARIAL REVIEW}}
\newcommand{\OPEN}{\taglabel{OPEN}}

\title{Observation Width Above Weight One:\\
Block Compression and the Residual Gap}
\author{a9lim}
\date{August 2026}

\begin{document}
\maketitle

\begin{abstract}
The flagship paper's transcript-span theorem proves that any
transcript-stable exact realizer of $\{Q=0\}$ observes refinement values
at directions spanning the input, so weight-$\leq w$ observations number
at least $\wt(x)/w$; the weighted-source \textsc{witt-fifo} rule attains
the bound at $w=1$, and the flagship claims no matching construction for
$w>1$.  This note proposes a resolution: for every fixed $w$, applying
the flagship realization theorem to a \emph{block-induced instance} --
partition the support of $x$ into blocks of weight at most $w$, take the
block indicator vectors as new coordinates, the Gram matrix of $B$ on
the blocks as the induced alternating form, and the block values of $Q$
as the induced diagonal -- yields a rule under the same
F1--F3/N1--N2 contract with $(w_0,c)=(w,1)$ whose distinct observations
are exactly the $\lceil\wt(x)/w\rceil$ block vectors.  Polarization makes
the induced diagonal reconstruct $Q(x)$ exactly, so exactness is
inherited, and the transcript-span bound is attained up to the integer
ceiling at every width.  The proof is a three-step reduction to the
flagship theorem and is recorded in full; it has not yet passed the
repository's adversarial review gate, and the residual questions --
all minor -- are listed at the end.
\end{abstract}

\section{The bound and its attainment at width one}

Notation and contract are the flagship's~\cite{goldarf}: $V=\F_2^n$ with
standard basis $e_1,\dots,e_n$, $B$ alternating, $Q$ a quadratic
refinement of $B$, input $x$ loaded by a refinement-blind map (F1), each
transition metered to at most $c$ oracle queries at Hamming weight at
most $w_0$ (F2), one declared outcome convention (F3), exactness uniform
over the refinement torsor (N1) within budget (N2).

\begin{theorem}[transcript-span lower bound]\label{thm:span}
\PROVED~\cite[transcript-span theorem]{goldarf}
Let the refinement family at fixed $B$ be closed under addition of every
linear functional.  If a transcript-stable rooted result is exact at
$x$, then $x\in\Span S(q,x)$; in the original frame, if every observed
vector has weight at most $w$, then $\wt(x)\leq|S(q,x)|\,w$.
\end{theorem}

\begin{theorem}[attainment at $w=1$]\label{thm:attain}
\PROVED~\cite[optimal observation support]{goldarf}
The weighted-source rule satisfies the contract at $(w_0,c)=(1,1)$ and
its distinct observations at input $x$ are exactly the $\wt(x)$
singleton directions $\{e_j:x_j=1\}$.
\end{theorem}

The flagship then records: for $w>1$ the theorem is only a lower bound;
no matching construction is claimed.  This note supplies the candidate
construction.

\section{The block-induced instance}

\begin{definition}[block partition]\label{def:blocks}
Fix $w\geq1$ and an input $x\in V$.  Partition $\supp(x)$ into blocks
$Z_1,\dots,Z_k$ of size at most $w$ by greedy chunking in coordinate
order, so $k=\lceil\wt(x)/w\rceil$.  Write $z_i=\sum_{j\in Z_i}e_j$ for
the block indicator vectors.  The partition is computed from $x$ alone.
\end{definition}

\begin{definition}[induced instance]\label{def:induced}
On $V'=\F_2^k$ with standard basis $e'_1,\dots,e'_k$, define
\[
 B'(e'_i,e'_j)=B(z_i,z_j)\quad(i\neq j),\qquad
 Q'\Bigl(\sum_{i\in s}e'_i\Bigr)=Q\Bigl(\sum_{i\in s}z_i\Bigr)
 \quad(s\subseteq\{1,\dots,k\}),
\]
and load the all-ones input $\mathbf{1}=e'_1+\dots+e'_k$.
\end{definition}

\begin{lemma}[the induced instance is well formed]\label{lem:induced}
\STANDARD\
$B'$ extends to an alternating bilinear form on $V'$, $Q'$ is a
quadratic refinement of $B'$ with diagonal $Q'(e'_i)=Q(z_i)$, and
\[
 Q'(\mathbf{1})=Q(x).
\]
\end{lemma}

\begin{proof}
$B'$ is the Gram matrix of $B$ on the $z_i$: it is symmetric, and its
diagonal $B(z_i,z_i)$ vanishes because $B$ is alternating, so $B'$ is
alternating.  The map $s\mapsto\sum_{i\in s}z_i$ is $\F_2$-linear and
injective (the $z_i$ have pairwise disjoint nonempty supports), so $Q'$
is the pullback of $Q$ along a linear embedding; pullbacks of quadratic
maps are quadratic, and the polar form of the pullback is the pullback
of the polar form, which is $B'$ by definition.  Explicitly, iterating
$Q(u+v)=Q(u)+Q(v)+B(u,v)$ gives
\[
 Q\Bigl(\sum_{i\in s}z_i\Bigr)
 =\sum_{i\in s}Q(z_i)+\sum_{\substack{i<j\\ i,j\in s}}B(z_i,z_j),
\]
which is the polarization expansion of the quadratic map on $V'$ with
diagonal $Q(z_i)$ and polar $B'$.  Since $z_1+\dots+z_k=x$, evaluating
at $s=\{1,\dots,k\}$ gives $Q'(\mathbf{1})=Q(x)$.
\end{proof}

\begin{proposition}[block compression]\label{prop:block}
\PROVED\ \UNREVIEWED\
Fix $w\geq1$.  Running the weighted-source \textsc{witt-fifo} rule on
the induced instance $(V',B',Q',\mathbf{1})$, with each induced
singleton oracle query $Q'(e'_i)$ implemented as the original-frame
query $Q(z_i)$, defines a rule satisfying F1--F3/N1--N2 at
$(w_0,c)=(w,1)$ whose root outcome is $P$ iff $Q(x)=0$, and whose
distinct refinement observations are exactly the $k=\lceil\wt(x)/w\rceil$
vectors $z_1,\dots,z_k$, each of original-frame weight at most $w$.
These observations have pairwise disjoint supports, are linearly
independent, and sum to $x$; the transcript-span bound of
Theorem~\ref{thm:span} is therefore attained up to the integer ceiling
at every width $w$.
\end{proposition}

\begin{proof}
\emph{Contract.}  F1: the block partition, the vectors $z_i$, the Gram
data $B'$, and the loaded input $\mathbf{1}$ are computed from $(B,x)$
and the fixed frame only; the flagship rule's own statics on the induced
instance consult only $(B',\mathbf{1})$ and are refinement-blind by the
flagship's locality proposition.  F2: the flagship rule queries its
oracle only at induced singletons $e'_i$, at most once per transition;
each such query is implemented as one original-frame query at $z_i$,
which has Hamming weight $|Z_i|\leq w$.  Thus $(w_0,c)=(w,1)$, constants
independent of $n$.  F3: the semantics is the flagship's single
normal-play claim convention, unchanged.  N1: for every refinement $Q$
of every $B$ and every $x$, Lemma~\ref{lem:induced} produces a
refinement $Q'$ of $B'$, and the flagship theorem is uniform over all
refinements of all alternating forms in all finite dimensions, so it
applies to $(V',B',Q',\mathbf{1})$; N2 holds with the stated budget.

\emph{Exactness.}  By the flagship realization theorem, the root of the
induced arena is $P$ iff $Q'(\mathbf{1})=0$, and
$Q'(\mathbf{1})=Q(x)$ by Lemma~\ref{lem:induced}.

\emph{Observation support.}  By the flagship's locality proposition, the
distinct observations of the induced run at input $\mathbf{1}$ are
exactly the induced singletons $\{e'_i:\mathbf{1}_i=1\}$, i.e.\ all of
them; under the implementation these are exactly $z_1,\dots,z_k$.
Disjoint supports give independence, and $\sum_iz_i=x$ by construction.
Finally $k=\lceil\wt(x)/w\rceil$, against the lower bound
$\wt(x)/w$.
\end{proof}

\begin{remark}[total observed weight]\label{rem:total}
The same construction is exact for the total-weight measure: the span
requirement forces $\sum_{z\in S}\wt(z)\geq\wt(x)$ for any exact
transcript-stable rule (the supports must cover $\supp(x)$), and the
block rule achieves $\sum_i\wt(z_i)=\wt(x)$ exactly.  So neither the
count-at-width nor the total-weight measure separates rules from
certificates once Proposition~\ref{prop:block} stands.
\end{remark}

\section{Status and residual questions}

Proposition~\ref{prop:block} is a new synthesis: its only inputs are the
flagship realization and locality theorems (proved, with their Lean-
checked ingredients) and the polarization expansion
(standard).  It has not yet been adversarially reviewed, and the
flagship should not absorb it before that gate.  Points a referee should
press:

\begin{enumerate}[leftmargin=*,itemsep=2pt]
\item that the contract's weight metering is genuinely frame-relative as
used here -- the flagship states its lower bound ``in the original
frame,'' and the block rule's queries are original-frame points of
weight $\leq w$, but the F2 clause should be re-read against the blocked
implementation;
\item that a per-input block partition is admissible under F1 -- the
loading map may consult $x$ and $B$ but not $Q$; greedy chunking
consults only $\supp(x)$;
\item that N1 uniformity survives the composition -- the induced family
$\{(B',Q')\}$ ranges over all alternating forms with all refinements as
$(B,Q,x)$ ranges, and the flagship theorem quantifies over exactly that
family.
\end{enumerate}

If the proposition survives, the flagship's second open question closes
positively and the residual gap is only the integer ceiling
$\lceil\wt(x)/w\rceil$ versus $\wt(x)/w$, which is not a real question.
Two sharper readings would remain open:

\begin{question}[uniform-in-width interface]\label{q:uniform}
\OPEN\ Is there a single rule, one arena family with one access
declaration, that attains the bound simultaneously at every width --
e.g.\ whose observation set at declared width $w$ is the block support
-- rather than one rule per fixed $w$?
\end{question}

\begin{question}[non-partition geometries]\label{q:geometry}
\OPEN\ The block rule observes a disjoint-support family.  Do
overlapping-support observation geometries buy anything -- for instance,
exactness with fewer than $\lceil\wt(x)/w\rceil$ queries on inputs with
special polar structure -- or does the span-plus-cover argument already
force the partition count on every input?  (For the uniform torsor
guarantee the lower bound applies verbatim; the question is only about
per-instance improvements below uniformity.)
\end{question}

This note is the authoritative statement of the width problem and will
be folded into the flagship paper -- as a corollary beside the $w=1$
attainment if Proposition~\ref{prop:block} survives review, as a
retracted attempt otherwise.

\begin{thebibliography}{9}\footnotesize
\bibitem{goldarf} a9lim, \emph{Quadratic refinements in normal play:
realization and the observation boundary}, companion flagship paper,
\path{writeups/goldarf.tex} in the ogdoad repository, 2026.
\bibitem{carlet2021} C.~Carlet, \emph{Boolean Functions for Cryptography
and Coding Theory}, Cambridge University Press, 2021.
\bibitem{zhang2010} Z.~Zhang, Y.~Shi, \emph{On the parity complexity
measures of Boolean functions}, arXiv:1004.0436 (2010).
\end{thebibliography}

\end{document}
